%% ma: L12-B19-0024
%% lop: 12
%% chuong: 6
%% bai: 19
%% dang: 12.19.10
%% loai: TLN
%% muc: VD
%% dap_an: "0,7"
%% nguon: "(NBV)-ĐỀ SỐ 15-MỖI NGÀY 1 ĐỀ THI 2025.pdf (D15, MỖI NGÀY 1 ĐỀ - Đề số 15), Phần III Câu 3"
%% kien_thuc: "Lớp 11 Bài 30 (nhân xác suất độc lập), Lớp 12 Bài 19 (công thức xác suất toàn phần)"
%% trang_thai: cho-duyet
%% ly_do: "Dạng mới 12.19.10 đề xuất (hai tên lửa độc lập bắn mục tiêu: xác suất toàn phần theo số quả trúng)."
%% ghi_chu: "Đáp án gốc 0,7 đúng. Bỏ ảnh máy bay phóng tên lửa (ảnh trang trí)."
\begin{ex}
Một máy bay bắn độc lập hai quả tên lửa vào một mục tiêu. Xác suất quả thứ nhất và quả thứ hai trúng đích lần lượt là $0{,}6$ và $0{,}7$. Nếu có một quả trúng đích thì mục tiêu bị tiêu diệt với xác suất $0{,}7$ và nếu có hai quả trúng đích thì xác suất này là $0{,}9$. Tính xác suất mục tiêu bị tiêu diệt.
\shortans{0,7}
\loigiai{Gọi $H_0,H_1,H_2$ lần lượt là các biến cố có $0$, $1$, $2$ quả trúng đích; $T$: ``mục tiêu bị tiêu diệt''.

$P(H_1)=0{,}6\cdot0{,}3+0{,}4\cdot0{,}7=0{,}46$; $P(H_2)=0{,}6\cdot0{,}7=0{,}42$; $P(T\mid H_0)=0$.

$P(T)=0{,}46\cdot0{,}7+0{,}42\cdot0{,}9=0{,}322+0{,}378=0{,}7$.}
\end{ex}
